\section{Background}
\label{sec:background}
New multichannel and digital offerings are, along with increased data availability, major developments that have characterized the retail sector in the past decades (\cite{Bradlow2017-qf,Cui2021-rw,Dekimpe2020-bh,Verhoef2010-wh}). Such advancements have given consumers access to a wide range of store formats and channels. In this thesis, \textit{store format} is a term referring to any type of physical store (e.g., hypermarket, supermarket, or convenience store) whereas \textit{channel} refers to the mode of access (e.g., online and offline). 

Retailers can operate through traditional physical store formats, through online-only channels, or in multichannel retailing settings (i.e., environments where online and offline channels operate in parallel but are not necessarily fully integrated, as in omnichannel retailing; \cite{Verhoef2015-mq}). For example, the grocery store is a traditional store format where consumers shop in person using carts in physical aisles. As an online-only format, Amazon’s early model was exclusively concerned with selling books online. A fashion brand such as H\&M could be taken as a multichannel retailer because its customers browse and buy items either through a website or in a physical store. As customers increasingly use multiple formats and channels, retailers gain more data to better understand customer decisions and improve performance evaluation across formats and channels alike.

With alternative store formats and channels, retailers can offer various forms of access for and engage a broader consumer base as they seek to enhance the shopping experience (\cite{Verhoef2015-mq}). As products become accessible through more physical locations and digital touchpoints, retailers also have the option of expanding store formats and adopting additional channels to increase the visibility of their offerings (\cite{Lemon2016-gc}). Therefore, retail formats and channels do not only provide access to products; they also determine how prominently those products are exposed to consumers during the shopping process. In particular, distinct retail formats and channels can facilitate different shopping goals, whether in the same consumers or among consumers whose characteristics and shopping behaviors differ from each other. It is thus beneficial for retailers to understand differences in how consumers behave across formats and channels, and to grasp these consumers differences in terms of their demographic factors and purchasing behaviors (\cite{Verhoef2010-wh}). Such insights lead to important implications for retailers who aim to effectively design and customize their marketing mixes with respect to their consumers’ needs (\cite{Breugelmans2023-fh}).

Retailers operating across store formats and channels have access to abundant data. This includes both macro-level data (e.g., total grocery purchases from particular demographic groups, or total number of global users on specific platforms) and micro-level data (e.g., individual or household purchases of certain categories, a single user’s consumption history). Such data-driven retailers generally have comprehensive information regarding products (what was bought), customers (who bought), channels (how it was bought), time (when it was bought), and locations (where it was bought from and delivered to). They may then use this data to understand their customer thoroughly, and to efficiently optimize their operations (\cite{Bradlow2017-qf,Van_Heerde2024-am}). Digital platforms such as Spotify or eBay also operate as retailers by offering products and services directly to customers, often without physical constraints. Platforms can leverage detailed customer data to produce better products and enhance customer loyalty (\cite{Smith2016-sp}). Similar to physical retailers’ decisions on shelf placement and in-store displays, digital platforms place products and content for users to experience through curated lists, rankings, and recommendations. In addition to data collected directly by retailers, publicly available information can be extracted from the internet via web scraping or APIs to generate further insights. These web-based data sources offer valuable context for understanding the market, consumer behavior, and competitive positioning across retailers (\cite{Bonfrer2022-tb,Guyt2024-nq}).

These developments offer opportunities and challenges for retailers and researchers alike. On one hand, diversity in retail formats allows businesses to tailor offerings for different customer segments and shopping contexts. For example, as \textcite{Forbes2024-nt} highlights, leading retailers are shifting from a one-size-fits-all model. As they customize store formats, product assortments, and services to meet the specific needs of segments such as Gen Z, millennials, and diverse cultural communities, retailers are becoming known as “support partners” for distinct groups of shoppers. On the other hand, in its volume and complexity, the available data demands a more structured, theory-driven approach to uncover meaningful insights. For example, \textcite{KPMG2025-ye} notes that leading retailers move beyond ad hoc or intuition-based decisions by developing structured data strategies aligned with business processes and decision-making. This approach ensures that data interpretation aligns with clear business objectives while reducing the influence of individual bias.

A deeper integration of insights is now required of consumer behavior research, empirical methods, and analytics as they focus on understanding consumer behavior, designing effective pricing strategies, and evaluating performance across channels. Insights act as a lens to interpret the observable patterns revealed by data, and data sharpens that focus by validating assumptions and uncovering causal relationships. Data thus becomes not just a tool for operational improvement but a means of systematic quantification regarding how customers respond to strategic decisions (and how those responses translate into retail performance).

\section{Problem Formulation}
Despite major advancements in data collection and marked expansion of retail formats and channels, using these developments to generate actionable insights about performance remains far from straightforward. Retailers still face several challenges in evaluating the outcomes of their strategic decisions. The challenges are especially pronounced when consumer behavior is used as the primary lens for evaluation. Because consumer behavior cannot always be directly observed, it must often be inferred from observable responses captured in data which often involves complex, context-specific responses that are difficult to isolate and interpret. Because consumer responses can vary widely across retail environments, retailers must go beyond tracking what happened to understanding why it happened, and eventually to grasping how it matters for future decisions (\cite{McKinsey-Company2025-qo}). Doing so requires an approach that integrates empirical analysis with established insights from consumer behavior literature to yield observations that are both theoretically sound and practically relevant. Four issues lend to the overall research problem which this thesis aims to address.

First, the growing reliance on data-driven approaches has not always been matched by the integration of theory and explanation. Retailers and researchers have access to vast volumes of behavioral data. However, while it may trace any number of behaviors—ranging from those exhibited during individual shopping experiences at physical stores to aggregate behavior on online platforms—more data does not inherently mean more insights. Without theoretical grounding, analysis can become reactive or descriptive at best. In other words, such analysis focuses on what happened, but not why. For example, a spike in sales following a price change might be observed, but without an understanding of reference price effects or shopper context, the implications for future pricing strategy remain unclear. As emphasized by \textcite{Bradlow2017-qf} and \textcite{Dekimpe2020-bh}, the challenge lies in combining empirical evidence with underlying mechanisms that explain why certain patterns emerge, when they matter, and how they generalize across contexts. In the absence of such integration, managerial decisions risk being guided by surface-level correlations rather than actionable guidance. Yet, to be actionable, guidance should be grounded in causal understanding, and it should be applicable across strategic settings.

Second, differences across store formats introduce challenges for both analysis and strategic execution. Even within the same retail chain, hypermarkets, supermarkets, and convenience stores cater to distinct shopping goals and customer expectations. For example, a customer may visit a hypermarket for planned, bulk shopping (e.g., detergent, milk) and a convenience store for quick, impulsive purchases (e.g., snacks). Despite the differences, retailers often apply uniform pricing and promotions across formats without fully accounting for how the same customer may respond differently depending on the setting. These challenges become even more pronounced when retailers target various customer segments by managing multi-tier private label (PL) portfolios that offer economy, standard, and premium products within the same category. Because the products are controlled by the retailer, pricing decisions across PLs tiers directly influence both customer choice and retailer margins. This creates uncertainty about the effectiveness of pricing strategies and raises a broader challenge. How should retailers design and evaluate pricing decisions when customer responses vary by format?

Third, when physical retailers introduce an online channel, it creates strategic uncertainty about long-term customer behaviors and business outcomes. While digital channels often increase convenience and expand access, they also influence how customers interact with the retailer over time. Some customers may shift their purchase volume online, while others may change product mix, basket size, and purchasing frequency (e.g. \cite{Ansari2008-qz,Li2015-tq,Melis2016-mf}). These behavioral shifts make it difficult to assess whether online channel initiatives effectively enhance long-term customer value or overall performance. The broader challenge is to isolate and quantify how such adoption changes purchasing behavior over time (in terms, e.g., of visit frequency, basket size, or category choice), and it impacts retailer revenue and profitability.

Fourth, consumer behavior on digital platforms is increasingly influenced by curation systems that determine which products, assortments, or contents receive exposure. In this context, exposure can be understood as a form of product visibility that shapes what consumers encounter in large digital assortments. For online-only environments like music streaming platforms, algorithmic and editorial curation shapes what users see and how they respond to products or content. Systems either amplify already-popular items or give visibility to emerging or independent content providers (e.g., \cite{Brynjolfsson2010-ya}). Platforms and retailers can thus act as gatekeepers of attention by influencing which products are noticed and ultimately consumed. Because consumer responses interact with such exposure, changes in the broader environment can also affect how consumers respond. External shocks such as the COVID-19 pandemic can shift user behavior and content consumption patterns. Accordingly, such shocks indirectly affect who gains or loses visibility. The broader challenge is to understand how, particularly in contexts where attention is scarce and competition is high, platform design and curator roles influence performance outcomes.

Together, these challenges point to a broader need to understand (by accounting for how consumers perceive, interpret, and respond) strategic retail decisions’ effects on business performance. Retailers influence consumer responses through multiple strategic instruments including pricing decisions, expansion of access through store formats and channels, and management of product visibility across physical and digital environments. This thesis thus responds to this broader need by drawing on established consumer behavior literature and applying empirical models to large-scale data to examine how pricing strategies, channel additions, and digital exposure influence performance outcomes such as sales and profitability. 

\section{Research Purpose and Questions}
The purpose of this thesis is to examine how strategic retail decisions, including pricing, channel introduction, and product exposure on digital platforms, influence customer behavior and thereby affect performance outcomes. By empirically modeling observable customer responses across store formats, channels, and digital platforms, the aim is to demonstrate how data-driven approaches, informed by consumer behavior research, can generate actionable insights for retailers.

This leads to the overarching research question: \textbf{How do retail strategies related to price, channel access, and product visibility shape customer behavior and ultimately influence retail performance across different formats, channels, and digital platforms?}

The thesis addresses this question through four empirical studies, each aligned with a specific sub-question:
\begin{enumerate}
    \item How do pricing strategies among PLs affect sales and profitability across different store formats within the same retail chain?
    \item How do regular price reductions and promotional discounts differentially affect sales outcomes across store formats within the same retail chain?
    \item What are the long-term effects of online channel adoption on customer behavior and retailer revenue?
    \item How did changes in product visibility and content on digital platforms affect customer behavior during the COVID-19 pandemic?
\end{enumerate}

\section{Outline of the Thesis}
The remainder of this thesis is organized as follows. Chapter~\ref{chap:theoryframework}, drawing on key concepts from consumer behavior research, introduces the theoretical framework to guide the empirical investigations. Chapter~\ref{chap:datamethod}, highlighting how different modeling approaches are applied in multiformat, multichannel, and digital platform contexts, outlines the data sources and empirical methods used across the four studies. Chapter~\ref{chap:study} presents the four studies and discusses their individual contributions in terms of the research sub-questions. Chapter~\ref{chap:discussion} jointly discusses the four studies’ theoretical contributions, managerial implications, methodological considerations, and limitations; it then notes their implications for future research. Finally, the Epilogue offers a personal reflection on the research process and a few confessions from crossing back and forth between academia and industry.